\documentclass[12pt]{article}
\usepackage{amsmath,amssymb,amsthm}
\usepackage[margin=1in]{geometry}

\newtheorem{theorem}{Theorem}
\newtheorem{lemma}[theorem]{Lemma}

\title{Problem 303: Multiples with Small Digits}
\author{Project Euler Solution}
\date{}

\begin{document}
\maketitle

\section{Problem Statement}
For a positive integer $n$, let $f(n)$ be the smallest positive multiple of $n$ whose decimal representation uses only the digits $\{0, 1, 2\}$.

Compute
\[
\sum_{n=1}^{10000} \frac{f(n)}{n}.
\]

\section{Mathematical Foundation}

\begin{theorem}[Existence]\label{thm:existence}
For every positive integer $n$, $f(n)$ exists and is finite.
\end{theorem}

\begin{proof}
The BFS algorithm below explores residue classes modulo~$n$, starting from residues corresponding to leading digits 1 and 2. Since there are only $n$ residue classes and each is visited at most once, the search terminates. The transition function $(10r + d) \bmod n$ for $d \in \{0,1,2\}$ generates a connected graph over $\mathbb{Z}/n\mathbb{Z}$ (since $\gcd(10, n)$ divides $10$ and the digit set $\{0,1,2\}$ spans all residues modulo any divisor of~10). Thus residue~0 is eventually reached, proving existence.
\end{proof}

\begin{lemma}[BFS correctness]\label{lem:bfs}
A breadth-first search building numbers digit-by-digit from $\{0,1,2\}$, tracking residues modulo~$n$, finds $f(n)$ correctly.
\end{lemma}

\begin{proof}
The state space is $\{0, 1, \ldots, n-1\}$. BFS starts with residues $1 \bmod n$ and $2 \bmod n$. At each step, from residue~$r$, we transition to $(10r + d) \bmod n$ for $d \in \{0,1,2\}$. BFS explores states in order of increasing digit count (and lexicographically within the same length). Each residue is visited at most once, so the first time residue~0 is reached yields the smallest qualifying multiple. The search visits at most $n$ residues.
\end{proof}

\begin{lemma}[Digit bound]\label{lem:digits}
For $n \le 10000$, $f(n)$ has at most $\sim 30$ digits. The hardest case is $n = 9999$, where $f(9999)$ has 28 digits.
\end{lemma}

\begin{proof}
Since $9999 = 3^2 \times 11 \times 101$, the smallest multiple of $9999$ using only digits $\{0,1,2\}$ is $f(9999) = \underbrace{1\cdots1}_{24}1112$, a 28-digit number. For all other $n \le 10000$, $f(n)$ has fewer digits, as verified computationally.
\end{proof}

\section{Algorithm}
\begin{verbatim}
function solve():
    total = 0
    for n = 1 to 10000:
        total += bfs_find_quotient(n)
    return total

function bfs_find_quotient(n):
    visited[0..n-1] = false
    queue = empty
    for d in {1, 2}:
        r = d mod n
        if r == 0: return d / n
        visited[r] = true; enqueue(r, bigint(d))
    while queue not empty:
        (r, value) = dequeue
        for d in {0, 1, 2}:
            nr = (10*r + d) mod n
            nv = 10*value + d
            if nr == 0: return nv / n
            if not visited[nr]:
                visited[nr] = true; enqueue(nr, nv)
\end{verbatim}

\section{Complexity Analysis}
\begin{itemize}
    \item \textbf{Time:} $O\!\left(\sum_{n=1}^{10000} n\right) = O(5 \times 10^7)$ residue transitions. Each transition costs $O(D)$ for big-integer updates where $D \le 30$.
    \item \textbf{Space:} $O(n)$ per query for the visited array and BFS queue.
\end{itemize}

\section{Answer}

\[
\boxed{1111981904675169}
\]

\end{document}
